\paragrafos{Δυνάμεις με εκθέτη φυσικό: όταν οι πολλαπλασιασμοί συνεχίζονται}{Α.7.8, «Δυνάμεις ρητών αριθμών με εκθέτη φυσικό»}

\begin{anagke}
Στο διάλειμμα η Πολυξένη λέει στην Αγγελική ένα μυστικό και της ζητάει να το πει μέσα σε μια ώρα σε δύο άλλα παιδιά. Κάθε παιδί που το μαθαίνει το λέει, μέσα στην επόμενη ώρα, σε δύο άλλα. «Μέσα σε δέκα ώρες θα το ξέρει όλο το σχολείο;» ρωτάει ο Ηλίας. Ο Στέλιος γράφει: $1,\ 2,\ 4,\ 8,\ 16,\ \ldots$ και κάθε φορά το πλήθος διπλασιάζεται. «Μετά από $10$ ώρες θα είναι $2\cdot2\cdot2\cdot\ldots$ δέκα φορές».

Ο Βασίλης λέει ότι $2$ δέκα φορές επί $2$ είναι $2\cdot10=20$. Ο Νίκος τον διορθώνει: «Όχι, δέκα φορές πολλαπλασιασμός, όχι δέκα φορές πρόσθεση. Θα βγει $1\,024$!». Η Ρήνη προτείνει μια σύντομη γραφή: $2\cdot2\cdot2\cdot2\cdot2=2^5$, το «δύο στην πέμπτη», που το ξέρουν από την Ε΄ και ΣΤ΄ Δημοτικού για τους φυσικούς. Η κυρία Αλεξάνδρα το βάζει στο νέο επίπεδο: «Μπορούμε να υψώσουμε στη δύναμη και έναν ρητό. Τι νομίζετε ότι είναι το $(-2)^3$; Και το $\left(-\frac{1}{2}\right)^3$;».

Σε αυτή την παράγραφο η παρέα ξαναβρίσκει τις δυνάμεις, μαθαίνει το πρόσημό τους και τις ιδιότητες που τις κάνουν γρήγορες, και προσέχει τα σημεία που μπερδεύουν τους περισσότερους: το $(-3)^2$ και το $-3^2$.
\end{anagke}

\begin{lexi}{δύναμη, βάση, εκθέτης}{δύναμη, βάση, εκθέτης}
Η \lexiword{δύναμη} είναι «ύψωση», δηλαδή το γινόμενο ίσων παραγόντων· στα αρχαία ελληνικά «δύναμις» ήταν και το τετράγωνο. Η \lexiword{βάση} είναι ο αριθμός που πολλαπλασιάζεται με τον εαυτό του και ο \lexiword{εκθέτης} (από το «εκ-τίθημι», «βάζω έξω») αυτός που γράφουμε πάνω δεξιά, δείχνοντας πόσους παράγοντες έχουμε.
\end{lexi}

\begin{metafrasi}
Πέντε φορές ο $2$ πολλαπλασιασμένος. & $2^5=2\cdot2\cdot2\cdot2\cdot2=32$\\
Ο $-2$ τρεις φορές. & $(-2)^3=(-2)(-2)(-2)=-8$\\
Μετά από ν γύρους το μυστικό το ξέρουν. & $2^\nu$ παιδιά\\
Πολλαπλασιάζω δυνάμεις με ίδια βάση. & $\alpha^\mu\cdot\alpha^\nu=\alpha^{\mu+\nu}$\\
Η δύναμη μιας δύναμης. & $(\alpha^\mu)^\nu=\alpha^{\mu\nu}$\\
\end{metafrasi}

\tmima{Η θεωρία}

\begin{orismos}{δύναμη με εκθέτη φυσικό}
Για κάθε ρητό $\alpha$ και φυσικό $\nu\ge2$ η \textbf{δύναμη} $\alpha^\nu$ είναι το γινόμενο $\nu$ παραγόντων ίσων με $\alpha$: $\alpha^\nu=\alpha\cdot\alpha\cdots\alpha$. Ορίζουμε ακόμη $\alpha^1=\alpha$.
\end{orismos}

Για παράδειγμα $(0{,}1)^3=0{,}1\cdot0{,}1\cdot0{,}1=0{,}001$ και $\left(\frac{2}{3}\right)^3=\frac{2}{3}\cdot\frac{2}{3}\cdot\frac{2}{3}=\frac{8}{27}$.

\tmimasmall{Το πρόσημο της δύναμης}

Από τον κανόνα των προσήμων του πολλαπλασιασμού:
\begin{itemize}[leftmargin=1.2em,itemsep=0.3ex,topsep=2pt]
\item Μια δύναμη με θετική βάση είναι πάντα θετική.
\item Μια δύναμη με αρνητική βάση είναι \textbf{θετική} αν ο εκθέτης είναι \textbf{άρτιος} και \textbf{αρνητική} αν είναι \textbf{περιττός}: $(-2)^4=16$, $(-2)^3=-8$, $(-1)^{100}=1$, $(-1)^{101}=-1$.
\end{itemize}

Η βάση είναι ό,τι βρίσκεται μέσα στην παρένθεση: $(-3)^2=9$ ενώ $-3^2=-(3\cdot3)=-9$.

\begin{idiotita}{ιδιότητες των δυνάμεων (εκθέτης φυσικός)}
Για ρητούς $\alpha,\,\beta$ και φυσικούς $\mu,\,\nu$ ισχύουν:
\begin{itemize}[leftmargin=1.2em,itemsep=0.3ex,topsep=2pt]
\item $\alpha^\mu\cdot\alpha^\nu=\alpha^{\mu+\nu}$ \qquad $\alpha^\mu:\alpha^\nu=\alpha^{\mu-\nu}$ (με $\alpha\neq0$, $\mu>\nu$)
\item $(\alpha\cdot\beta)^\nu=\alpha^\nu\cdot\beta^\nu$ \qquad $\left(\dfrac{\alpha}{\beta}\right)^\nu=\dfrac{\alpha^\nu}{\beta^\nu}$ (με $\beta\neq0$)
\item $(\alpha^\mu)^\nu=\alpha^{\mu\cdot\nu}$
\end{itemize}
\end{idiotita}

Γιατί; Το $2^3\cdot2^4=(2\cdot2\cdot2)\cdot(2\cdot2\cdot2\cdot2)=2^7$ (τρεις παράγοντες και τέσσερις: επτά). Το $(3^2)^3=3^2\cdot3^2\cdot3^2=3^6$. Το $(2\cdot5)^3=2\cdot5\cdot2\cdot5\cdot2\cdot5=2^3\cdot5^3=8\cdot125=1\,000$.

\tmimasmall{Προτεραιότητα πράξεων}

Η σειρά είναι: παρενθέσεις, μετά δυνάμεις, μετά πολλαπλασιασμοί και διαιρέσεις, τέλος προσθέσεις και αφαιρέσεις. Για παράδειγμα $2\cdot3^2-(-2)^3+10:(-5)=18-(-8)+(-2)=18+8-2=24$.

\begin{prosoxi}
\begin{itemize}[leftmargin=1.2em,itemsep=0.3ex,topsep=2pt]
\item $2^3=8$, όχι $6$: η δύναμη είναι επανάληψη πολλαπλασιασμού, όχι γινόμενο $2\cdot3$.
\item $(-3)^2=+9$, αλλά $-3^2=-9$ (ο εκθέτης ενεργεί μόνο στο $3$).
\item Το $(\alpha+\beta)^2$ δεν είναι $\alpha^2+\beta^2$: $(3+4)^2=49$ ενώ $3^2+4^2=25$.
\item Οι ιδιότητες $\alpha^\mu\cdot\alpha^\nu=\alpha^{\mu+\nu}$ ισχύουν μόνο για \emph{ίδια βάση}: το $2^3\cdot3^2$ δεν γίνεται $6^5$.
\item Στο $\alpha^\mu\cdot\alpha^\nu$ προσθέτω εκθέτες, δεν πολλαπλασιάζω: $2^3\cdot2^4=2^7$ και όχι $2^{12}$.
\end{itemize}
\end{prosoxi}

\tmima{Λυμένα παραδείγματα}

\paradeigma{Δυνάμεις και πρόσημα}
\tixero Οι δυνάμεις $2^5$, $(-3)^2$, $-3^2$, $(-3)^3$, $\left(-\frac{1}{2}\right)^3$, $(0{,}1)^3$.\par
\tithelo Να τις υπολογίσω.\par
\possindeo $2^5=32$. $(-3)^2=9$. $-3^2=-(3\cdot3)=-9$. $(-3)^3=-27$. $\left(-\frac{1}{2}\right)^3=-\frac{1}{8}$. $(0{,}1)^3=0{,}001$.\par
\apantisi $32$, $9$, $-9$, $-27$, $-\frac{1}{8}$, $0{,}001$.\par
\elegxos Αρνητική βάση με περιττό εκθέτη δίνει αρνητικό: το $(-3)^3$ και το $\left(-\frac{1}{2}\right)^3$.

\paradeigma{Με τις ιδιότητες}
\tixero Οι παραστάσεις $2^3\cdot2^4$, $5^8:5^5$, $(3^2)^3$, $(2\cdot5)^3$, $\left(\frac{2}{3}\right)^3$.\par
\tithelo Να τις υπολογίσω.\par
\possindeo $2^{3+4}=2^7=128$. $5^{8-5}=5^3=125$. $3^{2\cdot3}=3^6=729$. $2^3\cdot5^3=8\cdot125=1\,000$. $\frac{2^3}{3^3}=\frac{8}{27}$.\par
\apantisi $128$, $125$, $729$, $1\,000$, $\frac{8}{27}$.\par
\elegxos $(2\cdot5)^3=10^3=1\,000$.

\paradeigma{Το μυστικό του σχολείου}
\tixero Κάθε παιδί λέει το μυστικό σε $2$ νέα παιδιά κάθε ώρα. Στην αρχή το ξέρει $1$ παιδί.\par
\tithelo Πόσα παιδιά το ξέρουν μετά από $10$ ώρες.\par
\possindeo Μετά από $1$ ώρα: $2$. Μετά από $2$: $4$. Μετά από $3$: $8$. Μετά από $\nu$ ώρες το πλήθος διπλασιάζεται $\nu$ φορές: $2^\nu$. Για $\nu=10$: $2^{10}=1\,024$.\par
\apantisi $1\,024$ παιδιά.\par
\elegxos Από την $9η$ ώρα ($512$) στη $10η$ διπλασιάζεται: $1\,024$.

\paradeigma{Η προτεραιότητα}
\tixero Η παράσταση $2\cdot3^2-(-2)^3+10:(-5)$.\par
\tithelo Να την υπολογίσω.\par
\possindeo Πρώτα οι δυνάμεις: $3^2=9$, $(-2)^3=-8$. Μετά πολλαπλασιασμοί και διαιρέσεις: $2\cdot9=18$, $10:(-5)=-2$. Τέλος $18-(-8)+(-2)=18+8-2=24$.\par
\apantisi $24$.\par
\elegxos Αν έκανα πρώτα τον πολλαπλασιασμό $2\cdot3=6$, θα έβγαινε $6^2=36$: λάθος σειρά.

\tmima{Ασκήσεις}

\zoni{Το κατάλαβα}{δυνάμεις, πρόσημα, ιδιότητες}

\ex Υπολόγισε: α) $2^5$ \quad β) $3^4$ \quad γ) $(-2)^3$ \quad δ) $(-2)^4$ \quad ε) $(-1)^7$ \quad στ) $0^5$.

\ex Υπολόγισε: α) $(-3)^2$ \quad β) $-3^2$ \quad γ) $\left(-\frac{1}{2}\right)^3$ \quad δ) $(0{,}1)^3$ \quad ε) $(-0{,}2)^2$ \quad στ) $\left(\frac{2}{3}\right)^3$.

\ex Πες το πρόσημο χωρίς να υπολογίσεις: α) $(-5)^6$ \quad β) $(-5)^7$ \quad γ) $-(-3)^4$ \quad δ) $(-2)^3\cdot(-3)^2$.

\ex Γράψε ως μία δύναμη: α) $2^3\cdot2^4$ \quad β) $5^8:5^5$ \quad γ) $(3^2)^3$ \quad δ) $7^2\cdot7\cdot7^3$ \quad ε) $(-2)^5\cdot(-2)^2$ \quad στ) $(-4)^6:(-4)^4$.

\ex Υπολόγισε με τις ιδιότητες: α) $(2\cdot5)^3$ \quad β) $(-2\cdot3)^2$ \quad γ) $\left(\frac{2}{3}\right)^3$ \quad δ) $\left(-\frac{3}{4}\right)^2$.

\ex Υπολόγισε: α) $2\cdot3^2-(-2)^3+10:(-5)$ \quad β) $(-2)^2\cdot3-4^2:(-2)$ \quad γ) $5-3\cdot(-2)^2$ \quad δ) $\big((-1)^5+(-1)^6\big)\cdot10$.

\ex Σωστό ή λάθος; α) $2^3=6$. \quad β) $(-2)^4=16$. \quad γ) $-2^4=16$. \quad δ) $(\alpha\beta)^2=\alpha^2\beta^2$. \quad ε) $\alpha^2\cdot\alpha^3=\alpha^6$. \quad στ) $(\alpha^2)^3=\alpha^6$. \quad ζ) $(\alpha+\beta)^2=\alpha^2+\beta^2$.

\zoni{Το χρησιμοποιώ}{φήμες, κύβοι, σκάκι}

\ex \textbf{Το μυστικό.} Κάθε παιδί λέει το μυστικό σε $2$ νέα κάθε ώρα. Φτιάξε πίνακα με το πλήθος των παιδιών μετά από $0,\,1,\,\dots,\,10$ ώρες και πες πότε ξεπερνά τα $1\,000$.

\ex \textbf{Ο κύβος.} Ένας κύβος έχει ακμή $12$ cm. Βρες τον όγκο ($12^3$) και το εμβαδόν της επιφάνειάς του ($6\cdot12^2$).

\ex \textbf{Το σκάκι.} Στο πρώτο τετράγωνο ενός σκακιστικού πίνακα βάζουμε $1$ κόκκο, στο δεύτερο $2$, στο τρίτο $4$ και συνεχίζουμε διπλασιάζοντας. Πόσοι κόκκοι είναι στο $8$ο τετράγωνο και πόσοι στο $11$ο;

\ex \textbf{Τα βακτήρια.} Ένα βακτήριο τριπλασιάζεται κάθε ώρα. Πόσα είναι μετά από $5$ ώρες;

\ex \textbf{Οι τετράγωνοι αριθμοί.} Οι αριθμοί $1,\,4,\,9,\,16,\dots$ είναι τα τετράγωνα. Ποιος είναι ο $20$ός και ο $\nu$-οστός;

\ex \textbf{Η τιμή διπλασιάζεται.} Μια τιμή των $3$ € διπλασιάζεται $4$ φορές. Πόσο γίνεται;

\ex \textbf{Η δύναμη του $10$.} Υπολόγισε $10^3\cdot10^4$ και $10^8:10^5$.

\zoni{Το σκέφτομαι}{όταν το πρόσημο και ο εκθέτης παίζουν}

\ex Ο Βασίλης λέει ότι $2^3=6$. Βρες το λάθος και εξήγησε τη διαφορά ανάμεσα στο $2\cdot3$ και στο $2^3$.

\ex Δείξε με $\alpha=3$ και $\beta=4$ ότι $(\alpha+\beta)^2\neq\alpha^2+\beta^2$.

\ex Πότε ισχύει $(-\alpha)^\nu=\alpha^\nu$ και πότε $(-\alpha)^\nu=-\alpha^\nu$;

\ex Πόσο είναι τα $1^\nu$, $0^\nu$ (για $\nu\ge1$) και $(-1)^\nu$;

\ex Εξήγησε γιατί $2^3\cdot2^4=2^7$ γράφοντας όλους τους παράγοντες.

\ex Ποιο είναι μεγαλύτερο: το $3^4$ ή το $4^3$;

\zoni{Για το σπίτι}{δίπλωσε και μέτρα}

\ex Δίπλωσε ένα χαρτί στη μέση όσες φορές μπορείς. Πόσα στρώματα έχεις μετά από $1,\,2,\,3,\,4,\,5$ διπλώματα;
